\chapter{A projekt bemutatása}

A féléves projekt célja egy olyan \textbf{AES-alapú titkosító rendszer} elkészítése volt,
amely ugyanazt a funkciót több különböző végrehajtási úton tudja biztosítani:

\begin{itemize}
  \item \textbf{Natív, szekvenciális CPU implementáció} (C nyelven) -- ez adja az alap (baseline) viselkedést.
  \item \textbf{OpenCL-gyorsított megoldás} -- a párhuzamosítható részeket GPU-n futtatom.
  \item \textbf{Menedzselt .NET referencia} -- a beépített .NET kriptográfiai implementáció (összehasonlítási referencia).
\end{itemize}

A projekt két fő felhasználói felületet kínál: (1) egy \emph{single-file} titkosító/dekódoló részt,
amellyel fájlokat lehet csomagolt formában kezelni, valamint (2) egy \emph{benchmark} modult,
amely több iterációban méri a teljesítményt, és az eredményeket CSV-be menti.

\section{Kódbázis felépítése}

A repository több komponensből áll:

\begin{itemize}
  \item \textbf{\code{AES.WinForms}}: .NET WinForms alkalmazás, amely a GUI-t, a benchmark futtatást és a CSV export/importot kezeli.
  \item \textbf{\code{aeslib}}: natív C könyvtár az AES blokk-kódhoz és a módokhoz (CTR, GCM, CBC), valamint tesztek/benchmark/profiling.
  \item \textbf{\code{crypto\_aes\_opencl}}: OpenCL-es útvonal (host oldal C-ben + OpenCL kernel fájlok), külön tesztek és profiling.
  \item \textbf{\code{kernel\_loader}}: segédkönyvtár a kernel fájlok betöltéséhez (pl. futtatási környezettől függő útvonalkezelés).
  \item \textbf{\code{Docs}}: ez a dokumentum + a mért eredményeket tartalmazó CSV-k.
\end{itemize}

\section{Architektúra áttekintő ábra}

Az alábbi ábra a legfontosabb komponenseket és a hívási útvonalakat foglalja össze.

\begin{figure}[H]
  \centering
  \resizebox{0.94\textwidth}{!}{%
  \begin{tikzpicture}[
    node distance=12mm,
    box/.style={draw, rectangle, rounded corners=2pt, align=center, minimum width=4.2cm, minimum height=0.9cm},
    small/.style={draw, rectangle, rounded corners=2pt, align=center, minimum width=3.2cm, minimum height=0.8cm},
    arrow/.style={-Latex, thick}
  ]
    \node[box] (ui) {WinForms GUI\\\small \code{Form1}, szolgáltatások};
    \node[box, below=of ui] (bench) {BenchmarkService\\CSV export/import};
    \node[small, below left=of bench, xshift=-8mm] (managed) {Menedzselt\\\code{System.Security.Cryptography}};
    \node[small, below=of bench] (native) {Natív CPU\\\code{aeslib} (C)};
    \node[small, below right=of bench, xshift=8mm] (opencl) {OpenCL gyorsítás\\\code{crypto\_aes\_opencl}};

    \node[box, below=of opencl, yshift=4mm] (kernels) {OpenCL kernelfájlok\\\code{kernels/*.cl}};

    \draw[arrow] (ui) -- (bench);
    \draw[arrow] (bench) -- (managed);
    \draw[arrow] (bench) -- (native);
    \draw[arrow] (bench) -- (opencl);
    \draw[arrow] (opencl) -- (kernels);
  \end{tikzpicture}%
  }
  \caption{A projekt fő komponensei és a hívási irányok.}
\end{figure}

\section{A három kriptográfiai ``engine''}

A benchmark és a single-file modul ugyanazt a \emph{fogalmi} API-t használja: \emph{engine} + \emph{algoritmus} + \emph{paraméterek}.

\begin{description}
  \item[Natív CPU] Saját C implementációm, amely a logikát jól követhetően, \emph{szekvenciálisan} valósítja meg.
  \item[OpenCL] A titkosítás párhuzamosítható részét GPU-n futtatom. CTR-ben a blokk-függetlenség miatt ez különösen kézenfekvő.
        GCM-ben az AES-CTR rész és a GHASH feldolgozás egy része párhuzamosítható, viszont a tag előállítás/ellenőrzés miatt van szekvenciális komponens is.
  \item[Menedzselt AES] A .NET beépített implementáció (pl. \code{AesGcm}) modern CPU-n jellemzően AES-NI/specializált utasításokra támaszkodik,
        ezért ez \emph{nem baseline}, csak referencia.
\end{description}

\section{Mit mérek a benchmarkokban?}

Egy mérési \emph{session} során azonos paraméterek mellett több iterációban futtatok:

\begin{itemize}
  \item \textbf{Titkosítás} (\emph{Encrypt}) és \textbf{visszafejtés} (\emph{Decrypt})
  \item \textbf{Átviteli sebesség} (Throughput, \mbps)
  \item \textbf{Futási idő} (Elapsed, \ms)
  \item OpenCL esetén: \textbf{gyorsulás} a natív CPU baseline-hoz képest (\emph{SpeedupVsSequentialNativeCpu})
\end{itemize}

A dokumentum második fele ezekből a CSV-kből kiindulva elemzi az eltéréseket architektúra (x86/x64),
tápellátás (akkumulátor/töltő) és paraméterek (kulcsméret, adatméret, CTR vs GCM) szerint.

\cleardoublepage
